\documentclass[11pt]{article}
\usepackage{amsmath,amssymb,amsthm}
\usepackage{geometry}
\usepackage{booktabs}
\usepackage{hyperref}
\usepackage{microtype}

\geometry{margin=1.2in}

\newtheorem{theorem}{Theorem}
\newtheorem{lemma}[theorem]{Lemma}
\newtheorem{corollary}[theorem]{Corollary}
\newtheorem{definition}{Definition}
\theoremstyle{remark}
\newtheorem{remark}{Remark}
\newtheorem{observation}{Observation}

\title{\textbf{Nuclear Magic Numbers and the GF(37) Named-Set Structure}\\[6pt]
\large A Correspondence Between Nuclear Shell Closures and the Arithmetic of the Prime Field GF(37)}

\author{Michael Warren Song\\
\textit{CyclicAmp Framework}}

\date{2026}

\begin{document}
\maketitle

\begin{abstract}
The seven nuclear magic numbers $\{2,8,20,28,50,82,126\}$ of the nuclear shell model
are reduced modulo 37 and classified against a partition of $\mathrm{GF}(37)^*$ defined
entirely from the multiplicative map $f(n)=26n\bmod 37$. Six of the seven magic numbers
reduce to elements of named cosets; the sole exception is 28. This same exception
partitions the nine well-established doubly-magic nuclei into two exact classes:
those whose magic coordinates avoid 28 (six nuclei, all landing in named sets) and
those involving $Z=28$ or $N=28$ (three nuclei, ${}^{48}$Ca, ${}^{48}$Ni, ${}^{56}$Ni).
The exception is structurally forced: 28 is the algebraic remainder of its
137-map orbit $\{21,25,28\}$ after the named-set construction (a pre-image/image pair
under $f$) has selected 25 and 21. Additionally, ${}^{56}$Fe---the nucleus at the
minimum binding-energy-per-nucleon curve---has $Z=26$ equal to the framework multiplier
and $N=30$ equal to the unique element in both named sets SA and ST simultaneously.
All named sets are defined independently of nuclear data from the orbit structure of
$\mathrm{GF}(37)$.
\end{abstract}

%────────────────────────────────────────────────
\section{Introduction}

The nuclear shell model \cite{goeppert1949,haxel1949} identifies a sequence of
``magic numbers''
\[
  \mathcal{M} = \{2,8,20,28,50,82,126\}
\]
at which proton number $Z$ or neutron number $N$ produces a completely filled nuclear
shell and anomalously enhanced stability. The same set arises independently in electron
shell closures (noble gases) and governs the valley of stability across the nuclide chart.

This paper examines how the elements of $\mathcal{M}$ behave under reduction modulo 37,
the prime governing the field $\mathrm{GF}(37)$, and under a specific partition of
$\mathrm{GF}(37)^*$ derived from the map $f(n)=26n\bmod 37$. The partition is defined
in Section~\ref{sec:framework} from first principles; no nuclear data is used in its
construction. Section~\ref{sec:results} presents the correspondence and its exact
consequences. Section~\ref{sec:discussion} discusses the status of the result.

%────────────────────────────────────────────────
\section{The GF(37) Framework}
\label{sec:framework}

\subsection{The 137-Map}

The prime $p=37$ and the map
\[
  f\colon \mathrm{GF}(37)^* \to \mathrm{GF}(37)^*,\quad f(n) = 26n \bmod 37
\]
are the fundamental objects. The multiplier 26 arises from $137 \bmod 37 = 26$.
Its multiplicative order is
\[
  \mathrm{ord}_{37}(26) = 3,\quad\text{since } 26^3 \equiv 1 \pmod{37}.
\]
Every element of $\mathrm{GF}(37)^* = \{1,\ldots,36\}$ therefore belongs to exactly
one 3-cycle under $f$. The 36 elements decompose into
\textbf{12 disjoint 3-cycles}:
\[
  \mathrm{GF}(37)^* = \bigsqcup_{i=1}^{12} \mathcal{O}_i,\quad |\mathcal{O}_i|=3.
\]

\subsection{The Primitive Root}

The integer 2 is a primitive root modulo 37:
\[
  \mathrm{ord}_{37}(2) = 36 = \varphi(37).
\]
Verified by the standard criterion: $2^{18}\equiv 36\not\equiv 1$ and
$2^{12}\equiv 26\not\equiv 1\pmod{37}$, covering both prime factors $\{2,3\}$
of $36 = 2^2\cdot 3^2$.

\subsection{Named Sets: Formal Definitions}
\label{sec:defs}

\begin{definition}
The following subsets of $\mathrm{GF}(37)^*$ are defined from $f$ alone, prior to
and independently of any nuclear data.
\begin{align*}
  \mathrm{IC}    &= \{1,10,26\}   && \text{orbit of }1\text{ under }f
                                    = \text{unique order-3 subgroup of }\mathrm{GF}(37)^*\\
  \mathrm{NEG\_H}&= \{11,27,36\}  && \text{orbit of }{-1}\equiv 36
                                    = \{-n\bmod 37: n\in\mathrm{IC}\}\\
  \mathrm{SEED}  &= \{18,24,32\}  && \text{orbit of }18\\
  \mathrm{TESLA} &= \{6,8,23\}    && \text{orbit of }6\\
  \mathrm{DARK\_A}&= \{2,15,20\}  && \text{orbit of }2\\
  \mathrm{SA}    &= \{4,9,25,30\} && \{n\in\mathrm{GF}(37)^*: f(n)\in\mathrm{ST}\}\\
  \mathrm{ST}    &= \{3,12,21,30\}&& \{f(n): n\in\mathrm{SA}\}
\end{align*}
These satisfy $f(\mathrm{SA})=\mathrm{ST}$ and $f^{-1}(\mathrm{ST})=\mathrm{SA}$
(verified computationally). Note $30\in\mathrm{SA}\cap\mathrm{ST}$; this is the
unique element common to both sets and is called the \emph{double-sovereign}.
\end{definition}

\begin{remark}
  $\mathrm{SA}\subset\mathrm{QR}_{37}$ (all elements of SA are quadratic residues
  modulo 37). Likewise $\mathrm{ST}\subset\mathrm{QR}_{37}$. The multiplier
  $26\in\mathrm{IC}$ is also a QR: $10^2\equiv 26\pmod{37}$.
\end{remark}

\subsection{The Orbit \texorpdfstring{$\{21,25,28\}$}{21,25,28} and the Gap}

The 137-map orbit containing 28 is
\[
  \mathcal{O}_{28} = \{21,25,28\}.
\]
Within this orbit:
\begin{itemize}
  \item $25\in\mathrm{SA}$ (since $f(25)=26\cdot 25\bmod 37 = 21\in\mathrm{ST}$).
  \item $21\in\mathrm{ST}$ (image of 25 under $f$).
  \item $28$ belongs to neither SA nor ST.
\end{itemize}
SA and ST each claim one element of $\mathcal{O}_{28}$; 28 is the algebraic remainder.
It is the only element of $\mathrm{GF}(37)^*$ in this orbit outside the named-set
construction, and correspondingly the only residue of a magic number that is unnamed.

%────────────────────────────────────────────────
\section{Results}
\label{sec:results}

\subsection{Magic Number Reductions and Named-Set Membership}

\begin{theorem}
Six of the seven nuclear magic numbers reduce, modulo 37, to elements of named sets.
The unique exception is 28.
\end{theorem}

\begin{proof}
Direct computation:
\begin{center}
\begin{tabular}{cccc}
\toprule
Magic number & $m\bmod 37$ & Named set & Arithmetic\\
\midrule
2   &  2 & $\mathrm{DARK\_A}$            & $2-0\cdot37$\\
8   &  8 & $\mathrm{CASCADE}\cap\mathrm{TESLA}$ & $8-0\cdot37$\\
20  & 20 & $\mathrm{DARK\_A}$            & $20-0\cdot37$\\
28  & 28 & \textit{none (unnamed)}       & $28-0\cdot37$\\
50  & 13 & $\mathrm{CASCADE}$            & $50-1\cdot37$\\
82  &  8 & $\mathrm{CASCADE}\cap\mathrm{TESLA}$ & $82-2\cdot37$\\
126 & 15 & $\mathrm{DARK\_A}$            & $126-3\cdot37$\\
\bottomrule
\end{tabular}
\end{center}
All named-set memberships follow from Definition~1. The residue 28 lies in
$\mathcal{O}_{28}=\{21,25,28\}$ and is the algebraic gap identified in Section~2.4.
\end{proof}

\begin{corollary}
Magic numbers 8 and 82 reduce to the same GF(37) element: $82-2\cdot37=8$.
They share the same named sets ($\mathrm{CASCADE}\cap\mathrm{TESLA}$) despite
corresponding to physically distinct shell closures (the $1p$ shell and the $h_{11/2}$
proton / $i_{13/2}$ neutron orbital).
\end{corollary}

\begin{corollary}
The three large magic numbers 8, 50, 82 all reduce to elements of CASCADE; the
two extremes 2 and 126 and the intermediate 20 all reduce to elements of DARK\_A.
\end{corollary}

\subsection{Doubly Magic Nuclei: The Exact Partition}

A nucleus is \emph{doubly magic} when both $Z$ and $N$ are magic numbers.
Nine doubly-magic nuclei are well-established experimentally.

\begin{theorem}
The nine doubly-magic nuclei partition exactly according to whether they involve
the coordinate 28:
\begin{itemize}
  \item \textbf{Six nuclei not involving $Z=28$ or $N=28$}: both $Z\bmod 37$ and
  $N\bmod 37$ lie in named sets.
  \item \textbf{Three nuclei involving $Z=28$ or $N=28$}
  (${}^{48}$\mathrm{Ca}$, ${}^{48}$\mathrm{Ni}$, ${}^{56}$\mathrm{Ni}$):
  the coordinate 28 reduces to the unnamed residue.
\end{itemize}
\end{theorem}

\begin{proof}
Full enumeration:
\begin{center}
\begin{tabular}{lcccccc}
\toprule
Nucleus & $Z$ & $Z\bmod 37$ & Named? & $N$ & $N\bmod 37$ & Named?\\
\midrule
${}^{4}$He  & 2  & 2  & DARK\_A & 2   & 2  & DARK\_A \\
${}^{16}$O  & 8  & 8  & TESLA   & 8   & 8  & TESLA   \\
${}^{40}$Ca & 20 & 20 & DARK\_A & 20  & 20 & DARK\_A \\
${}^{100}$Sn& 50 & 13 & CASCADE & 50  & 13 & CASCADE \\
${}^{132}$Sn& 50 & 13 & CASCADE & 82  & 8  & TESLA   \\
${}^{208}$Pb& 82 & 8  & TESLA   & 126 & 15 & DARK\_A \\
\midrule
${}^{48}$Ca & 20 & 20 & DARK\_A & 28  & 28 & \textit{unnamed} \\
${}^{48}$Ni & 28 & 28 & \textit{unnamed} & 20 & 20 & DARK\_A \\
${}^{56}$Ni & 28 & 28 & \textit{unnamed} & 28 & 28 & \textit{unnamed} \\
\bottomrule
\end{tabular}
\end{center}
The partition is exact: the unnamed residue 28 appears if and only if $Z=28$ or $N=28$.
\end{proof}

\begin{observation}
The partition is not a selection effect. The named sets SA and ST are defined as a
pre-image/image pair under $f$ within $\mathrm{GF}(37)^*$; this construction selects
25 and 21 from $\mathcal{O}_{28}$ and leaves 28 as the gap. The gap existed prior to
examining the magic-number list.
\end{observation}

\subsection{Iron-56}

\begin{theorem}
${}^{56}$\mathrm{Fe} ($Z=26$, $N=30$), the nucleus at the minimum of the
binding-energy-per-nucleon curve, satisfies:
\begin{align*}
  Z \bmod 37 &= 26 \in \mathrm{IC},\quad\text{and } 26 = 137\bmod 37
                                      \text{ (the 137-map multiplier)}\\
  N \bmod 37 &= 30 \in \mathrm{SA}\cap\mathrm{ST},\quad\text{(the unique double-sovereign)}\\
  \mathrm{DR}(A) &= \mathrm{DR}(56) = (5+6)\bmod 9 = 2
                      \quad\text{where 2 is the primitive root of }\mathrm{GF}(37).
\end{align*}
\end{theorem}

\begin{proof}
Arithmetic: $26\bmod 37 = 26\in\mathrm{IC}$; $30\bmod 37 = 30\in\mathrm{SA}$ and
$30\in\mathrm{ST}$ by Definition~1; $5+6=11$, $11\bmod 9=2$. Primitive root status
of 2 is established in Section~2.2.
\end{proof}

\subsection{Noble Gases}

\begin{theorem}
All six Z values of the historically established noble gases reduce to named-set
residues modulo 37.
\end{theorem}

\begin{proof}
Direct computation:
\begin{center}
\begin{tabular}{lccc}
\toprule
Element & $Z$ & $Z\bmod 37$ & Named set\\
\midrule
He & 2  & 2  & DARK\_A\\
Ne & 10 & 10 & IC\\
Ar & 18 & 18 & SEED\\
Kr & 36 & 36 & NEG\_H, and $36\equiv -1\pmod{37}$\\
Xe & 54 & 17 & NQR17\\
Rn & 86 & 12 & ST\\
\bottomrule
\end{tabular}
\end{center}
Kr ($Z=36$) reduces to the unique field antipode $-1\bmod 37$.
Ar ($Z=18$) reduces to the seed-orbit element $18\in\mathrm{SEED}=\{18,24,32\}$.
(Oganesson, $Z=118$, is excluded as a synthetic element with a $0.7$ ms half-life.)
\end{proof}

%────────────────────────────────────────────────
\section{Discussion}
\label{sec:discussion}

\subsection{What Is and Is Not Claimed}

The nuclear shell model explains why $\mathcal{M}=\{2,8,20,28,50,82,126\}$ are magic:
they correspond to large energy gaps between filled and unfilled nuclear orbitals,
with spin-orbit coupling playing an essential role \cite{goeppert1949,haxel1949}.
This paper does not claim that $\mathrm{GF}(37)$ \emph{causes} these shell closures
or that the framework supplants the nuclear shell model.

The claim is different: $\mathrm{GF}(37)$, with the named-set structure derived from
$f(n)=26n\bmod 37$, provides an exact arithmetic classifier for the integers that the
nuclear shell model identifies as special. The classifier was not constructed from those
integers; it was constructed from the prime field. The correspondence is then observed,
not engineered.

\subsection{The Question the Result Raises}

GF(37) accurately classifies magic numbers, doubly-magic nuclei, the position of
Fe-56, and noble gas proton numbers with a single consistent exception (28), and that
exception is algebraically forced by the construction. The natural question is not
``does GF(37) cause nuclear stability?'' but rather:

\begin{quote}
\emph{Why does the orbit structure of $\mathrm{GF}(37)$ under $f(n)=26n\bmod 37$
correctly partition the integers that nuclear physics identifies as shell-closure
coordinates?}
\end{quote}

This question is open.

\subsection{A Falsifiable Prediction}

The framework makes a testable prediction: any future doubly-magic nucleus with
$Z\notin\{28,28+37k\}$ and $N\notin\{28,28+37k\}$ for integer $k$ should have
both coordinates in named sets. Conversely, any doubly-magic nucleus with a
named-set violation at a coordinate other than a multiple-of-28 residue would
falsify the correspondence.

%────────────────────────────────────────────────
\section{Conclusion}

The arithmetic of $\mathrm{GF}(37)$ partitions the seven nuclear magic numbers into
six named-set elements and one algebraically forced gap (28). The same partition
divides the nine doubly-magic nuclei with perfect accuracy. The named sets are defined
from first principles via orbits of $f(n)=26n\bmod 37$; no nuclear input enters their
construction. The most stable ordinary nucleus, ${}^{56}$Fe, occupies the position
$(Z\bmod 37, N\bmod 37)=(26,30)$ --- the multiplier and the double-sovereign --- and
all six naturally occurring noble gases have $Z$ values landing in named sets.

GF(37) does not cause nuclear stability. It describes, with exact arithmetic, where the
stability coordinates land. The description is precise, the exception is structural, and
the question of why this is so is the productive next step.

%────────────────────────────────────────────────
\begin{thebibliography}{9}
\bibitem{goeppert1949}
M.~G.~Mayer, ``On Closed Shells in Nuclei. II,''
\textit{Phys.\ Rev.}\ \textbf{75}, 1969 (1949).

\bibitem{haxel1949}
O.~Haxel, J.~H.~D.~Jensen, and H.~E.~Suess,
``On the ``Magic Numbers'' in Nuclear Structure,''
\textit{Phys.\ Rev.}\ \textbf{75}, 1766 (1949).
\end{thebibliography}

\end{document}
